\section{Correctness and Index Privacy}
\label{sec:security}
\subsection{Conditional Correctness}
\label{sec:corr}
Let \(e_i\in R\) be the actual independently sampled encryption errors.
Before any coefficient reduction modulo \(q\), define in the integer ring \(R\)
\begin{equation}
\label{eq:integer-lifts}
 \widehat{PM}_j=\sum_{\nu=1}^{\alpha}B^{\nu-1}\widehat P_{i_{\nu},j},\qquad
 \widehat E_j^{\rm agg}=\sum_{i=1}^N\widehat P_{i,j}e_i,\qquad
 Z_j=\widehat{PM}_j+t\widehat E_j^{\rm agg}.
\end{equation}
These products reduce only by \(X^n+1\). Under the radix capacity condition,
\(\widehat{PM}_j\) is also the canonical lift of \(PM_j\): every coefficient
lies in \([0,B^\alpha)\), hence
\(\|\widehat{PM}_j\|_\infty\le B^\alpha-1\). Noise products are negacyclic convolutions
and can have negative integer coefficients.

\begin{theorem}[Conditional correctness]
\label{thm:correctness}
Fix admissible \(pp\), a legal database and target tuple, and any realization
of key generation and query encryption. Assume that for every \(j\in[L]\),
\begin{equation}
\label{eq:no-wrap}
 \|\widehat{PM}_j\|_\infty+t\|\widehat E_j^{\rm agg}\|_\infty<q/2.
\end{equation}
A sufficient, data-independent plaintext bound is
\[
 B^\alpha-1+t\|\widehat E_j^{\rm agg}\|_\infty<q/2
 \qquad(j\in[L]).
\]
Under \eqref{eq:no-wrap}, Algorithms~\ref{alg:qgen}--\ref{alg:replyext} return
\((m_{i_1},\ldots,m_{i_\alpha})\).
\end{theorem}
\begin{proof}
For each query ciphertext, \(c_{0,i}+c_{1,i}s\equiv v_{I,i}+te_i\pmod{qR}\).
If \(\mathsf{Resp}_j=(d_{0,j},d_{1,j})\), the ring arithmetic in
Algorithm~\ref{alg:replygen} gives
\begin{align*}
 d_{0,j}+d_{1,j}s
 &\equiv\sum_i\widehat P_{i,j}(v_{I,i}+te_i)\pmod{qR}\\
 &\equiv\widehat{PM}_j+t\widehat E_j^{\rm agg}=Z_j\pmod{qR}.
\end{align*}
This is a congruence; it does not identify the integer linear combination
of canonical ciphertext representatives with \(Z_j\).
By \eqref{eq:no-wrap}, \(\|Z_j\|_\infty<q/2\). Thus
\(\operatorname{center}_q(d_{0,j}+d_{1,j}s)=Z_j\) in \(R\), and
\[
 \mathsf{Dec}_s(\mathsf{Resp}_j)=[Z_j]_t
    =[\widehat{PM}_j]_t=PM_j.
\]
For each \(u\), its canonical coefficient is exactly
\(y_{j,u}=\sum_{\nu}k_{\nu}x_{i_{\nu},(j-1)n+u+1}\), without plaintext wraparound.
Lemma~\ref{lemma:packing-condition} recovers all target segments, including
zero padding. The Encode/Decode identity then recovers every target record.
\end{proof}

Direct encryption of \(k_{\nu}\) does not multiply its error \(e_i\) by \(k_{\nu}\).
The weights increase the plaintext term; the data polynomials determine
noise accumulation through \(\sum_i\widehat P_{i,j}e_i\), including errors
from non-target coordinates. Bounding a value only after centering would
not rule out prior wraparound, and is not the argument above.

\subsection{Parametric Full-Task Correctness}
\label{sec:parametric-correctness}
We now bound the failure of one complete ordered \(\alpha\)-record
retrieval under the formal \(\chi_r\), for any \(r>0\). Theorem~\ref{thm:correctness}
remains the deterministic implication used below; no numerical \(r\)
or cryptanalytic security estimate is needed for this analysis.

\begin{lemma}[Rounded-noise tails]
\label{lem:rounded-tail}
Let \(Y\sim\mathcal N(0,r^2)\) and \(Z=\operatorname{Round}(Y)\).
For every integer \(h\ge0\),
\begin{equation}
\label{eq:rounded-strict-tail}
 p_h(r):=\Pr(|Z|>h)
 =\operatorname{erfc}\!\left(\frac{h+1/2}{\sqrt2\,r}\right).
\end{equation}
For integer \(h\ge1\), \(\Pr(|Z|\ge h)=p_{h-1}(r)\), with threshold
\((h-1/2)/(\sqrt2\,r)\) in the complementary error function.
\end{lemma}
The proof is given in Appendix~\ref{app:lem-rounded-tail}.


Write \(e_i=\sum_{u=0}^{n-1}e_{i,u}X^u\) and define the single event
\begin{equation}
\label{eq:global-good-event}
 G_h=\{\,|e_{i,u}|\le h\text{ for every }i\in[N],\ 0\le u<n\,\}.
\end{equation}
Coefficient-wise and draw-wise independence give
\begin{equation}
\label{eq:global-event-probability}
 \Pr(G_h^c)=1-(1-p_h(r))^{Nn}\le Nn\,p_h(r).
\end{equation}
All \(N\) query errors, including non-target coordinates, enter this event.
The same errors affect every reply block; neither \(L\) nor \(\alpha\)
introduces an additional union factor.

\begin{lemma}[Block support and negacyclic noise]
\label{lem:block-noise}
For \(j\in[L]\), let
\[
 m_j=\min\{n,\max\{0,J-(j-1)n\}\}.
\]
Then \(1\le m_j\le n\) and
\(\|\widehat P_{i,j}\|_1\le m_j(B-1)\).
For any signed \(P,e\in\mathbb Z[X]/(X^n+1)\),
\(\|Pe\|_\infty\le\|P\|_1\|e\|_\infty\).
Consequently, on \(G_h\), simultaneously for all \(j\),
\begin{equation}
\label{eq:aggregate-good-event-bound}
 \|\widehat E_j^{\rm agg}\|_\infty\le N m_j(B-1)h.
\end{equation}
\end{lemma}
The proof is given in Appendix~\ref{app:lem-block-noise}.


\begin{proposition}[Full-task correctness under rounded Gaussian errors]
\label{prop:full-task-correctness}
Fix admissible \(pp\), a legal database \(M\), and a legal ordered tuple
\(I=(i_1,\ldots,i_\alpha)\). Use the formal iid law \(\chi_r\), \(r>0\),
and generate a fresh key, query and reply by the stated algorithms.
If integer \(h\ge0\) satisfies
\begin{equation}
\label{eq:sufficient-h}
 B^\alpha-1+tN m_j(B-1)h<q/2\qquad(j\in[L]),
\end{equation}
then the probability that \(\mathsf{ReplyExt}(pp,sk,\mathsf{Resp},\mathsf{st})\)
fails to return \((m_{i_1},\ldots,m_{i_\alpha})\), in that order, is at most
\[
 1-(1-p_h(r))^{Nn}.
\]
The probability includes KeyGen, all query-encryption errors, all \(a_i\),
and all other formal encryption randomness.
\end{proposition}
\begin{proof}
On \(G_h\), Lemma~\ref{lem:block-noise} and
\(\|\widehat{PM}_j\|_\infty\le B^\alpha-1\) make
\eqref{eq:sufficient-h} imply \eqref{eq:no-wrap} for every block.
Theorem~\ref{thm:correctness} then recovers every packed block and every
requested record in the prescribed order. Task failure is thus contained
in \(G_h^c\), whose probability is \eqref{eq:global-event-probability}.
This implication holds for every realization of the secret and the
uniform \(a_i\); averaging over them adds no failure term.
\end{proof}

\paragraph{Exact integer threshold.}
Set \(A=B^\alpha-1\), \(D_j=tN m_j(B-1)>0\), and
\(C=(q-1)/2-A\). Since \(t=2^w\) and \(\gcd(t,q)=1\) make \(q\) odd,
\[
 A+D_jh<q/2\quad\Longleftrightarrow\quad A+D_jh\le(q-1)/2.
\]
If \(C<0\), no nonnegative integer \(h\) passes this sufficient screen;
this does not prove impossibility. Otherwise define
\begin{equation}
\label{eq:hmax}
 h_{\max,j}=\left\lfloor\frac{C}{D_j}\right\rfloor,\qquad
 h_{\max}=\min_{j\in[L]}h_{\max,j}.
\end{equation}
Then \(h_{\max}\) passes every block and \(h_{\max}+1\) fails at least
one limiting block. Because \(p_h(r)\) decreases with \(h\), this choice
gives the strongest bound within this uniform coefficient-event method.
Define the resulting ideal-law failure upper bound
\begin{equation}
\label{eq:epsilon-dec}
 \epsilon_{\rm dec}(r)
 =1-\left(1-\operatorname{erfc}
 \left(\frac{h_{\max}+1/2}{\sqrt2\,r}\right)\right)^{Nn}.
\end{equation}
It is the exact probability of \(G_{h_{\max}}^c\), not necessarily the
actual decryption-failure probability.

\begin{corollary}[Union bound and a sufficient noise region]
\label{cor:correctness-region}
Under the proposition's hypotheses,
\(\Pr[\mathrm{task\ failure}]\le Nn\,p_h(r)\).
When \(C\ge0\), for any ideal target \(0<\varepsilon<1\), the region
\begin{equation}
\label{eq:r-correctness-region}
 0<r\le
 \frac{h_{\max}+1/2}
 {\sqrt2\,\operatorname{erfc}^{-1}(\varepsilon/(Nn))}
\end{equation}
suffices for \(\epsilon_{\rm dec}(r)\le\varepsilon\).
\end{corollary}
\begin{proof}
Use \eqref{eq:global-event-probability}. On \((0,1)\),
\(\operatorname{erfc}^{-1}\) is positive and decreasing, so
\eqref{eq:r-correctness-region} is equivalent to
\(Nn\,p_{h_{\max}}(r)\le\varepsilon\).
\end{proof}

The exact-product bound also gives a symbolic endpoint: put
\(\delta_*=1-(1-\varepsilon)^{1/(Nn)}\) and replace
\(\varepsilon/(Nn)\) by \(\delta_*\) in
\eqref{eq:r-correctness-region}. This is equivalent to
\(\epsilon_{\rm dec}(r)\le\varepsilon\) for the bound function and is
at least as permissive as the union region. No active-profile numerical
inverse or noise value is selected here.


\subsection{Normal Form and Ciphertext Pseudorandomness}
\label{sec:normal-form}
Let \(\mathcal F_{\rm real}^{(\tau)}=((A_i,A_i s+E_i))_{i=1}^{\tau}\)
have \(s\leftarrow\chi\), independent \(E_i\leftarrow\chi\), and independent
uniform \(A_i\in R_q\). Let \(\mathcal F_{\rm unif}^{(\tau)}\) be uniform
in \((R_q^2)^\tau\), and define \(\Adv^{\rm NF}_{\mathcal D}(\kappa;\tau)\)
as the corresponding acceptance-probability gap. This is an intermediate
distribution, not an additional hardness assumption.
Write
\[
 p_{\rm unit}:=\Pr_{a\leftarrow U(R_q)}[a\in R_q^\times].
\]
The normal-form method is established in the discussion preceding
Definition~1 of~\cite{homomorphic} and in
\cite[Lemma~2.24 and Claim~2.25]{lpr2013toolkit}.
The following discrete-polynomial specialization records its sample count
and loss in our notation; it is not a new reduction.

\begin{corollary}[Normal-form specialization]
\label{cor:normal-form}
For every PPT \(\mathcal D\) on \(\tau\) normal-form samples, there is a
PPT \(\mathcal B\) on \(\tau+1\) uniform-secret samples such that
\[
 \Adv^{\rm NF}_{\mathcal D}(\kappa;\tau)
 \le \frac{\Adv^{\rm uPLWE}_{\mathcal B}(\kappa;\tau+1)}
              {p_{\rm unit}}.
\]
\end{corollary}
See Appendix~\ref{app:normal-form} for the specialization proof.

This consumes one additional \emph{ring} sample. Lemma~2.24 of
\cite{lpr2013toolkit} is stated with continuous errors and a discretized
secret; it does not automatically make both laws the same discrete
\(\chi\). The transformation in Appendix~\ref{app:normal-form} verifies that property for our
already-discrete input samples.
For cyclotomic \(R\), Claim~2.25 of~\cite{lpr2013toolkit} supplies
\(p_{\rm unit}\ge 1/\operatorname{poly}(n,\log q)\) for any integer
\(q\ge2\). For the special case of prime \(q\equiv1\pmod{2n}\) and
\(f=X^n+1\), it gives \(p_{\rm unit}=(1-1/q)^n\ge e^{-1}\).
We use the general lower bound for the construction; prime \(q\) is not
a theorem requirement. For a general monic \(f\), the quantitative bound
still retains \(p_{\rm unit}\), without asserting this cyclotomic lower bound.

\begin{corollary}[XPIR-SKE ciphertext pseudorandomness]
\label{cor:ske-pseudorandom}
If \(\gcd(t,q)=1\), pseudorandomness of the normal-form joint distribution
implies XPIR-SKE ciphertext pseudorandomness for any fixed message vector
in \(R_t^\tau\).
The transformation preserves the distinguisher's advantage and sample count.
\end{corollary}
The bijective scaling argument is in Appendix~\ref{app:ske}.

The message vector may also be fixed conditional on public, sample-independent
side information, as in Definition~\ref{def:mq-cpir-privacy}.

\subsection{Reduction from Uniform-secret PLWE}
\begin{theorem}[Single-challenge index privacy of the formal construction]
\label{thm:mq-privacy}
Fix the formal construction with \(f=X^n+1\), \(n\) a power of two,
a specified efficiently sampleable \(\chi\), and \(\gcd(t,q)=1\).
For every PPT adversary \(\mathcal A\) in
Definition~\ref{def:mq-cpir-privacy}, there is a PPT \(\mathcal B\)
using \(N+1\) uniform-secret PLWE samples such that
\[
 \Adv^{\rm priv}_{\Pi,\mathcal A}(\kappa)
   \le \frac{\Adv^{\rm uPLWE}_{\mathcal B}(\kappa;N+1)}
                  {p_{\rm unit}}.
\]
Here \(N,n,\log q\), the public parameter descriptions, and the encoded
database length are polynomially bounded in \(\kappa\).
Consequently, the corresponding \((N+1)\)-sample uniform-secret
assumption in Definition~\ref{def:uplwe} implies single-challenge index
privacy, since \(p_{\rm unit}\) has the inverse-polynomial lower bound
stated in Section~\ref{sec:normal-form}.
\end{theorem}

\begin{proof}[Proof outline]
The normal-form reduction provides \(N\) shared-secret pairs from
\(N+1\) uniform-secret samples. Scaling by the unit \(t\) and translating
by the selected constants simulates the entire query jointly, with no
per-target hybrid. A uniform query and its deterministic reply are
independent of the challenge bit. The normal-form abort branch multiplies
both acceptance probabilities by \(p_{\rm unit}\), giving the stated
loss under the two declared advantage conventions. Appendix~\ref{app:privacy}
gives the complete simulation, including illegal inputs, the unconditional
uniform acceptance probability, and polynomial-time conditions.
\end{proof}


The dependency chain is uniform-secret PLWE with \(N+1\) samples,
the existing normal-form reduction to \(N\) shared-secret samples,
the SKE scaling and message-translation step, and finally the reduction
to packed index privacy above. The extra ring sample is a reduction
input, not an additional ciphertext sent by the client.

\paragraph{Scope.}
The proof establishes only the defined single-challenge index privacy under
the corresponding uniform-secret assumption for \((f,q,\chi)\).
It makes no worst-case ideal-lattice reduction claim and does not require
prime \(q\). It does not establish adaptive
multi-round security, CCA security, SPIR, malicious correctness, or a QROM
claim. Unpack, Decode, and all client decryption results are absent from
the adversary's view. Correctness, no-wrap, and decryption-failure probability
therefore play no role in this privacy proof. Theorem~\ref{thm:correctness}
remains a separate conditional correctness result.

\paragraph{Key lifetime.}
Each PIR query uses a fresh secret key, authorized only for that query's
\(N\) SKE ciphertexts. Completion, abort, or retry retires its encryption
authority; a retry or next query runs KeyGen again. No additional diagnostic
or public encryptions under that key are permitted. Thus a completed query
has \(\tau_{\rm visible}=N\) ciphertexts per key (an aborted query has at most
\(N\)); Theorem~\ref{thm:mq-privacy} still uses \(N+1\) uPLWE input samples.
Retaining the key to decrypt its outstanding reply does not authorize
further encryption. This policy introduces no key-reuse or multi-query
security claim.
